%% Revision log for hecke-algebras-SL2bar.tex (Draft 5, November 2026).
%% Cross-references to the paper use xr. Compile the paper first.
\documentclass[12pt,letterpaper]{amsart}
\usepackage[T1]{fontenc}
\usepackage{mathptmx}
\DeclareMathAlphabet{\mathcal}{OMS}{cmsy}{m}{n}
\usepackage[letterpaper,margin=1in]{geometry}
\usepackage{amsmath,amssymb}

\usepackage{enumitem}
\usepackage[final]{microtype}
\usepackage{booktabs}
\usepackage{longtable}
\usepackage{xr}
\externaldocument[P-]{hecke-algebras-SL2bar}

\newcommand{\Gbar}{\overline{G}}
\newcommand{\fbar}{\overline{f}}
\newcommand{\cO}{\mathcal{O}}
\DeclareMathOperator{\pr}{pr}
\newcommand{\cH}{\mathcal{H}}

\begin{document}

\title{Revision log for ``On Hecke algebras of the $p$-adic metaplectic
cover $\overline{SL_2}$'', Draft 5}
\author{David Joyner and Claude}
\date{November 24, 2026}
\maketitle

This log accompanies Draft~5 of the paper (file
\texttt{hecke-algebras-SL2bar.tex}). The paper is a corrected and
expanded version of D.~Joyner, \emph{A correspondence for the
generalized Hecke algebra of the metaplectic cover
$\overline{SL(2,F)}$, $F$ $p$-adic}, New York J. Math.\ \textbf{4}
(1998), 223--235, cited below as [Jo98]. The working source of [Jo98]
was the file \texttt{howe6.tex}. Numbers such as
``Theorem~\ref{P-thrm:main}'' refer to Draft~5 of the paper.
Section~\ref{sec:conc34} gives the correspondence with the numbering of
Drafts~2 to~4.

\section{Drafts}

\begin{description}[leftmargin=1.5em,style=nextline]
\item[Original] [Jo98], from \texttt{howe6.tex}.
\item[Draft 1] (November 23, 2026). File
\texttt{hecke-alg-metaplectic-SL2-corrected.tex}. Mathematical
correction in response to the referee report by Carbon. The main
theorem of [Jo98] was replaced. The revision log was an appendix of the
paper.
\item[Draft 2] (November 23, 2026). File
\texttt{hecke-algebras-SL2bar.tex}. Revision to the house style guide
(Draft~2 of November 11, 2026), new title, Claude added as co-author,
one proof corrected, literature review added. This log is now a
separate file.
\item[Draft 3] (November 23, 2026). Same file. Response to Carbon's
second referee report, which reviewed Draft~2. Glossary rewritten,
Question~\ref{P-q:n2} sharpened, several proofs clarified. The
footnote and table of contents changes made after Draft~2 was first
posted are also counted here.
\item[Draft 4] (November 24, 2026). Same file. Restructured into two
parts at David Joyner's request: Part~I for $n=2$ along the lines of
[Jo98], Part~II for $n\geq 3$. New results for $n=2$ and a new
non-existence theorem for every $n\geq 2$. Changes requested by David
Joyner to Sections~6.3, 8 and~9 of Draft~3.
\item[Draft 5] (November 24, 2026). Same file. At David Joyner's
request, the paper now states its motivation and explains why the
comparison with $SL(2)$ fails and how far off it is. New:
Table~\ref{P-tab:howfar}, Proposition~\ref{P-prop:traceless} and three
remarks.
\end{description}

\noindent
A different revision of [Jo98], made earlier in 2026 with Argon and
ChatGPT, is in the folder \texttt{hecke-alg-corresp} of the repository
\texttt{metaplectic-jl}. Drafts~1--5 here were made from
\texttt{howe6.tex} and do not incorporate that revision.

\section{AI usage}

\begin{itemize}[leftmargin=1.5em]
\item \textbf{Gemini (Google), ``Carbon''.} Refereed [Jo98] for
correctness and readability. The report recommended rejection or major
revision. It is answered point by point in
Section~\ref{sec:referee}. Carbon also refereed Draft~2. That report
recommended acceptance subject to minor revisions, and it is answered in
Section~\ref{sec:referee2}.
\item \textbf{Claude Fable 5.2 (Anthropic).} Produced Draft~1. It
checked the referee report, decided with David Joyner to replace the
main theorem by what survives, wrote the new statements and proofs of
Sections~\ref{P-section:genuine}--\ref{P-section:savin}, and rewrote
Section~\ref{P-section:sl2}.
\item \textbf{Claude Opus 5.5 (Anthropic).} Produced Draft~2. It
re-derived every proof in the paper (verification level~1), found and
corrected one error in a proof (item~\ref{it:coverb} of
Section~\ref{sec:d1d2}), wrote the literature review, background
examples, notation table, summary and glossary, checked bibliographic
data by web search, and wrote this log. It also produced Draft~3,
checking each point of Carbon's second report before acting on it.
It produced Drafts~4 and~5 as well, including the new results of
Section~\ref{P-section:n2} and Proposition~\ref{P-prop:traceless}.
\item \textbf{Carbon} also supplied a written summary of [Sa88] for
Draft~4. It was used only for orientation (Section~\ref{sec:sources}).
\item \textbf{David Joyner} chose the route of an honest corrected
draft over local fixes, chose the new title, and asked that Claude be
listed as an author.
\end{itemize}

\noindent
Verification levels follow the style guide. Level~1 is an independent
derivation. Level~2 is a check against a source actually consulted.
Level~3 is reliance on memory. No independent human reader has checked
the new arguments.

\section{Original and new results}

\subsection*{Results of [Jo98] retained, with corrections}
\begin{itemize}[leftmargin=1.5em]
\item Lemma~\ref{P-lemma:general}, the Howe--Moy multiplicativity
lemma (proof added).
\item Lemma~\ref{P-lemma:subgroups}(a), (b), (d), corrected and with
proofs.
\item Proposition~\ref{P-prop:main}. The generators and relations
A(i), A(ii), B, C(i) are retained. A(i), C(ii), C(iii) are corrected.
C(iv) is replaced.
\item Proposition~\ref{P-prop:borel}(a)--(c), from Borel. Part (d) is
corrected.
\item The statement of Savin's theorem, corrected as
Theorem~\ref{P-thm:savin}.
\end{itemize}

\subsection*{Results new in 2026}
Lemma~\ref{P-lemma:hilbert}, Lemma~\ref{P-lemma:cover},
Lemma~\ref{P-lemma:subgroups}(c), (e), Lemma~\ref{P-lemma:length},
Example~\ref{P-ex:t0}, Lemma~\ref{P-lemma:canonical},
Definition~\ref{P-def:admissible}, Lemma~\ref{P-lemma:support},
Lemma~\ref{P-lemma:genuine} (a corrected form of [Jo98, Lemma~3]),
Proposition~\ref{P-prop:admissible}, Proposition~\ref{P-prop:main2} (it
replaces [Jo98, Proposition~2]), Theorem~\ref{P-thrm:main},
Corollary~\ref{P-corollary:main}, Example~\ref{P-ex:trace},
Proposition~\ref{P-prop:level0}, and
Questions~\ref{P-q:defining}, \ref{P-q:n2}, \ref{P-q:savin},
\ref{P-q:reps}. New in Draft~4: Lemma~\ref{P-lemma:nonsplit},
Proposition~\ref{P-prop:Ksub}, Lemma~\ref{P-lemma:signs},
Theorem~\ref{P-thm:unimodular} and Corollary~\ref{P-cor:unimodular}.
New in Draft~5: Proposition~\ref{P-prop:traceless}.

\section{Changes of status}

\begin{enumerate}[leftmargin=1.5em]
\item \emph{[Jo98, Theorem~1].} The map $f_g\mapsto\fbar_g$ was
claimed to be an algebra bijection whenever $\overline{B_k}$ splits.
Withdrawn. It is false for $n\geq 3$
(Theorem~\ref{P-thrm:main}, Corollary~\ref{P-corollary:main}). The case
$n=2$ is open (Question~\ref{P-q:n2}).
\item \emph{[Jo98, Corollary~1].} Bijection of irreducible
representations with equal traces. Withdrawn. False for $n\geq 3$
(Example~\ref{P-ex:trace}).
\item \emph{[Jo98, Corollary~2].} Preservation of supercuspidal and
square-integrable representations. Withdrawn, since it depended on the
false theorem and its proof confused vectors with matrix coefficients.
Now an open question (Question~\ref{P-q:reps}).
\item \emph{[Jo98, Proposition~1], final claim.} That the relations
form a defining set. Downgraded to an open question
(Question~\ref{P-q:defining}).
\item \emph{[Jo98, Proposition~4(d)].} Irreducibility of the induced
module. Replaced by the statement about its unique irreducible
quotient.
\item \emph{[Jo98, Corollary~3]} (traces at Iwahori level). Removed.
\item \emph{Draft~3, Question~10.1} ($n=2$). Partly answered in
Draft~4. For constants of absolute value one the answer is negative,
for every $n\geq 2$ (Theorem~\ref{P-thm:unimodular}). In particular
the main theorem of [Jo98] also fails for $n=2$. The case of other
constants, and abstract isomorphism, remain open
(Question~\ref{P-q:n2}).
\end{enumerate}

\section{Concordance}

Old numbers are those of the printed [Jo98], which numbers each kind
of statement separately. These are the numbers used in the referee
report.

\begin{longtable}{@{}p{0.3\textwidth}p{0.62\textwidth}@{}}
\toprule
{[Jo98]} & Draft 2 \\
\midrule
\endhead
Definition 1 & Definition~\ref{P-def:1} \\
Definition 2 & Definition~\ref{P-def:genuineHecke} \\
Lemma 1 & Lemma~\ref{P-lemma:subgroups}(b) \\
Lemma 2 & Lemma~\ref{P-lemma:general}(a), (d) \\
Lemma 3 & Lemma~\ref{P-lemma:genuine} \\
Lemma 4 & Definition~\ref{P-def:genuineHecke} and the last formula of
Lemma~\ref{P-lemma:genuine} \\
Lemma 5 & Remark~\ref{P-rem:cocycle} \\
Proposition 1 & Proposition~\ref{P-prop:main} \\
Remark 1 & Remark after Example~\ref{P-ex:t0}. Item (5) removed \\
Lemma 6 & Lemma~\ref{P-lemma:subgroups}(a), (d) and
Lemma~\ref{P-lemma:general}(b), (c) \\
Proposition 2 & Proposition~\ref{P-prop:main2} \\
Theorem 1 & Withdrawn. See Proposition~\ref{P-prop:admissible} and
Theorem~\ref{P-thrm:main} \\
Corollary 1 & Withdrawn. See Corollary~\ref{P-corollary:main} \\
Definition 3, Proposition 3 & Removed (matrix coefficients) \\
Definition 4 & Text before Proposition~\ref{P-prop:borel} \\
Proposition 4 & Proposition~\ref{P-prop:borel} \\
Definition 5, Corollary 2 & Removed. See Example~\ref{P-ex:trace} and
Question~\ref{P-q:reps} \\
Theorem 2 & Theorem~\ref{P-thm:savin} \\
Corollary 3 & Removed \\
\bottomrule
\end{longtable}

\section{Response to Carbon's first referee report (on [Jo98])}
\label{sec:referee}

Each item gives the report's claim, our verdict (correct, misprint,
fixable gap, or false), the verification level, and the action taken.

\subsection*{Major mathematical errors}
\begin{enumerate}[leftmargin=1.5em]
\item \emph{The main theorem is false.} Correct for $n\geq 3$, level~1.
The torus commutator $(a,b)_n^2$ is Lemma~\ref{P-lemma:cover}(a). The
report's suggestion that $n=2$ also fails is not established. Action:
Theorem~\ref{P-thrm:main}, Corollary~\ref{P-corollary:main},
Question~\ref{P-q:n2}.
\item \emph{Savin's theorem is misstated.} Correct, level~2 (via the
summary of Savin's result in Karasiewicz's paper, not via Savin's
paper). Action: Theorem~\ref{P-thm:savin} and
Remark~\ref{P-vrem:savin}.
\item \emph{At Iwahori level $\fbar_{t_0}$ is unsupported when $n$ is
even.} False for $n=2$, since then $(u,\varpi)_2^2=1$. True for every
$n\geq 3$, odd or even, level~1. Action:
Proposition~\ref{P-prop:level0}.
\item \emph{At positive level $B_kt_0B_k$ may fail to support a genuine
function.} The criterion stated is correct (it is admissibility,
Definition~\ref{P-def:admissible}). The feared failure does not occur
at deep level, level~1. Action: Proposition~\ref{P-prop:admissible}.
\item \emph{The generators $\fbar_g$ are not well defined or not
bi-invariant.} Correct, level~1. Action: the functions
$\fbar_{\tilde g}$ of Lemma~\ref{P-lemma:support}.
\item \emph{Lemma~3 ignores the zero function and the cocycle
factor.} Correct, level~1. Action: Lemma~\ref{P-lemma:genuine}, whose
last formula records the factor $\beta(g_1,g_2)^{-1}$.
\item \emph{The span of the $\fbar_x$, $x\in B$, is a twisted group
algebra.} Correct, level~1. Action: Proposition~\ref{P-prop:main2},
relation ($\overline{\mathrm{B}}$)(ii).
\item \emph{Proposition~2, A(ii): stray $g$ in a domain of
integration.} Misprint. The proof containing it is replaced.
\item \emph{The lift of $-I$ squares to $(-1,-1)_n$.} Correct, level~1.
This is a cocycle phase on the subalgebra over $B$, covered by
($\overline{\mathrm{B}}$)(ii). It is not by itself an obstruction.
\end{enumerate}

\subsection*{Additional issues}
\begin{enumerate}[leftmargin=1.5em,resume]
\item \emph{Remark~1(5) and relation C(iv) are false for $k\geq 1$.}
Correct, level~1. Action: Example~\ref{P-ex:t0} and the new C(iv).
\item \emph{The normal form $\tilde f_w=f_{s_1}^a*f_{t_0}^b$ is wrong.}
Correct, level~1. Action: factorizations in the proof of
Proposition~\ref{P-prop:main}.
\item \emph{A(i) constant: $q$ versus $c$.} Correct, the constant is
$q$, level~1. Also the unbound index $i$ in C(ii), C(iii) was a
misprint.
\item \emph{Lemma~1: the displayed map is identically zero.} Misprint,
corrected in Lemma~\ref{P-lemma:subgroups}(b).
\item \emph{Corollary~1: surjectivity not proved, $\phi$ written for
$\eta$.} Correct but moot, since the corollary is withdrawn.
\item \emph{Corollary~2 confuses vectors with matrix coefficients.}
Correct, level~1. The corollary is withdrawn.
\item \emph{Proposition~4(d): $I(W)$ need not be irreducible.} Correct,
level~3 for the general theory. Action:
Proposition~\ref{P-prop:borel}(d).
\item \emph{Remarks on compactness of the center.} Correct but not
needed. Not used.
\end{enumerate}

\subsection*{Readability and minor corrections}
All adopted. The uniformizer is now $\varpi$ and the projection is
$\pr$. The Haar measures are normalized once. The undefined symbol
$\overline{G_s}$ is gone. The typographical errors listed in the report
(``anonomous'', ``Iwawori'', ``muliplicity'', ``contruction'',
``denotes'', ``follow'', ``implies'', ``are also are'', ``This
established'', the doubled ``the'', the missing parentheses in $x(g)$,
\texttt{W\^{}a}, ``Bruhat decomposition'') are corrected or removed with
the text that contained them. The document is converted from
\texttt{nyj2e} to \texttt{amsart}.

\subsection*{The cocycle formula}
The report first questioned the formula for $\beta$ and then verified
it. We agree with the conclusion. Level~1 for the diagonal case
$\beta(h(a),h(b))=(b,a)_n$ and for $U^\pm$
(Lemma~\ref{P-lemma:cover}). The general cocycle identity is cited from
Kubota (Remark~\ref{P-vrem:kubota}).

\subsection*{Recommendation}
The report recommended rejection or a complete rethinking of the main
theorem. We adopted the second option. We do not adopt the report's
unqualified statement that $\cH(G/\!/B_k)$ and the genuine algebra are
not isomorphic. What is proved is the weaker statement of
Corollary~\ref{P-corollary:main} about support-preserving maps.

\section{Response to Carbon's second referee report (on Draft 2)}
\label{sec:referee2}

The report recommended acceptance subject to minor revisions. It
stated that every numbered proof had been checked and found sound. Each
point is answered below with a verdict and a verification level.

\subsection*{Technical points}
\begin{enumerate}[leftmargin=1.5em]
\item \emph{Hypothesis of Savin's theorem.} The summary calls this a
slip, but the body of the report concludes that $\gcd(p,n)=1$ is the
right hypothesis. Verdict: no error found. We add one observation
(level~1). When $n$ is even, $\gcd(p,n)=1$ already forces $p$ to be
odd, so the hypotheses $p\nmid n$ and $p\nmid 2n$ differ only when $n$
is odd and $p=2$. The report quotes the first page of [Sa88], but it
does not say that the source was consulted, so we treat the quotation
as level~3. Remark~\ref{P-vrem:savin} is kept.
\item \emph{Question~\ref{P-q:n2} ($n=2$).} The report suggests
distinguishing maps given by a choice of lifts, whose scalars lie in
$\mu_2$, from maps with arbitrary non-zero scalars. Adopted. The
question now asks about arbitrary scalars, and it gives the example of
the lift of $-1$ over $\mathbb{Q}_2$, whose square is
$(-1,-1)_2=-1$. The computation of that square from the cocycle is
level~1. The report's discussion of Gauss sums at Iwahori level was not
checked and is not used.
\item \emph{Lemma~\ref{P-lemma:canonical}, case $p\mid n$.} The report
re-derives the argument and finds it correct, including for $g$ outside
$GL(2,\cO)$. It suggests stating why $\log$ and $\exp$ commute with
conjugation. Adopted (level~1). The paragraph before
Definition~\ref{P-def:deep} now gives the continuity argument and
separates the conditions for $\exp$ and for $\log$.
\item \emph{Remarks~\ref{P-vrem:kubota} and \ref{P-vrem:borel}.} The
report states that the cited results and section numbers agree with
[Ku67], [Bo76] and [BH06]. The report does not say that these sources
were consulted, and the style guide asks for claims of another AI to be
checked independently. A partial check: the table of contents of [BH06]
has a section titled ``The Hecke Algebra'' in its first chapter,
consistent with the citation of \S4.3, but the title of \S4.3 itself
was not seen. Both remarks are kept.
\end{enumerate}

\subsection*{Glossary}
The report found most entries too short and proposed replacements.
Adopted in substance. Every entry now has two to four sentences, and an
entry for Kubota's cocycle was added. The proposed text was not used
verbatim, for four reasons. It pointed to numbered statements, which
the house rules exclude from glossary entries. It used semicolons. One
sentence in the entry for genuine functions was garbled. The entry for
types asserted that the trivial representation of $B_k$ is not a
single-block type in general, which we have not checked, so the new
entry makes no claim about $B_k$.

\subsection*{Readability and minor items}
\begin{enumerate}[leftmargin=1.5em]
\item \emph{Claude on the author line.} The report notes that arXiv and
most journals do not accept AI authors. Not changed, since this was
David Joyner's decision. The acknowledgments already describe the AI
contributions. A version for submission would need the change.
\item \emph{Proof of Lemma~\ref{P-lemma:cover}(b).} The step
$(a,-1)_n^{-2}=1$ is now written as
$(a,-1)_n^{-2}=(a,(-1)^{-2})_n=(a,1)_n=1$. Adopted.
\item \emph{Proof of Lemma~\ref{P-lemma:subgroups}(e).} Added that
$B_k$ and $uB_ku^{-1}$ both contain $\Gamma_{k+1}$ because $u\in K$
normalizes $\Gamma_{k+1}$. Adopted (level~1).
\item \emph{Order of $s_1$.} The comparison with [HM85] now refers to
the permutation matrix used there for $GL(n)$. Adopted.
\item \emph{Subtitle.} The report reminds us to update the subtitle when
finalizing. The draft number is now~3. The rest of the subtitle follows
the style guide and is unchanged.
\item \emph{[Wa24].} Kept as an arXiv citation. We did not check
whether it has since been published.
\item \emph{Proof of Lemma~\ref{P-lemma:hilbert}.} The report approves
it. No change.
\end{enumerate}

\section{Concordance from Draft 3 to Draft 5}
\label{sec:conc34}

Numbers not listed are unchanged. Draft~4 differs from Draft~5 only
in Example~8.8 (now Example~\ref{P-ex:unimodular}), Theorem~11.3 (now
Theorem~\ref{P-thm:savin}) and Remark~11.4 (now
Remark~\ref{P-rem:savin}). The right-hand column gives the Draft~5
numbers.

\begin{longtable}{@{}p{0.3\textwidth}p{0.62\textwidth}@{}}
\toprule
{Draft 3} & Draft 5 \\
\midrule
\endhead
Lemma 3.5 & Lemma~\ref{P-lemma:subgroups} \\
Lemma 3.6 & Lemma~\ref{P-lemma:length} \\
Example 3.7 & Example~\ref{P-ex:subgroups} \\
Definition 3.8 & Definition~\ref{P-def:1} \\
Example 3.9 & Example~\ref{P-ex:hecke} \\
Lemma 3.10 & Lemma~\ref{P-lemma:general} \\
Theorem 7.3 & Theorem~\ref{P-thrm:main} \\
Corollary 7.4 & Corollary~\ref{P-corollary:main} \\
Example 7.5 & Example~\ref{P-ex:q7} \\
Section 8, Proposition 8.1 & Section~\ref{P-section:reps},
Proposition~\ref{P-prop:borel} \\
Example 8.3 & Example~\ref{P-ex:trace} \\
Section 9, Proposition 9.1 & Section~\ref{P-section:savin},
Proposition~\ref{P-prop:level0} \\
Example 9.2 & Example~\ref{P-ex:level0} \\
Theorem 9.3 & Theorem~\ref{P-thm:savin}, with its last sentence moved
to Remark~\ref{P-rem:savin} \\
Section 10 & Section~\ref{P-section:summary} \\
Questions 10.1--10.3 & Questions~\ref{P-q:n2}, \ref{P-q:savin},
\ref{P-q:reps} \\
\bottomrule
\end{longtable}

\section{Response to David Joyner's comments on Draft 4}

David Joyner asked that the paper state its motivation. The genuine
theory of the cover was known to resemble that of $PGL(2)$, or more
generally of $G/Z_n$, rather than that of $SL(2)$. The 1998 paper
asserted a correspondence with $SL(2)$ at the level of Hecke algebras,
and the question is why that fails and how far off it is. He asked for
this to be explained in the introduction and in other places.

\begin{enumerate}[leftmargin=1.5em]
\item The introduction has a new subsection ``Motivation'', which
states the 1998 assertion and that it is false, and a new subsection
``How far off is the comparison with $SL(2)$?'' with
Table~\ref{P-tab:howfar}.
\item Remark~\ref{P-rem:twist} describes the genuine algebra as the
linear one twisted by Kubota's cocycle, seen through double cosets.
\item Remark~\ref{P-rem:n2where} explains where the obstruction lives
for $n=2$.
\item Proposition~\ref{P-prop:traceless} is new. For $n\nmid 2m$ the
Hecke operators attached to $t_0^m$ have trace zero in every admissible
genuine representation. It generalizes Example~\ref{P-ex:trace}
(level~1).
\item Remark~\ref{P-rem:samelattice} relates this to the supports at
Iwahori level and to Savin's theorem. The translations $t_0^m$ with
$n\nmid 2m$ are missing at Iwahori level and traceless at deep level.
\item The abstract and the summary are updated accordingly.
\end{enumerate}

\noindent
Three statements were checked and weakened before release. The table
entry for Iwahori-level supports when $n=2$ says only that no
translation is excluded, because the paper proves a necessary condition
for support, not a sufficient one. For the same reason the abstract and
Remark~\ref{P-rem:samelattice} do not claim that every $t_0^m$ with
$n\mid 2m$ is seen by traces. In Remark~\ref{P-rem:n2where}, the element
$w(\varpi)$ is $-t_0s_1$, not $t_0s_1$.

\section{Changes from Draft 4 to Draft 5}

\begin{enumerate}[leftmargin=1.5em]
\item New subsections ``Motivation'' and ``How far off is the
comparison with $SL(2)$?'' in the introduction, with
Table~\ref{P-tab:howfar}.
\item New Remarks~\ref{P-rem:twist}, \ref{P-rem:n2where} and
\ref{P-rem:samelattice}, new Proposition~\ref{P-prop:traceless} with a
following remark.
\item The abstract and the summary of results are extended.
\item Numbering changes are listed in Section~\ref{sec:conc34}.
\end{enumerate}

\section{Response to David Joyner's comments on Draft 3}

\begin{enumerate}[leftmargin=1.5em]
\item \emph{Restrict first to $n=2$ and prove what can be proven along
the lines of the original paper.} Done as Part~I. The general
constructions of Sections~\ref{P-section:sl2}--\ref{P-section:metaplectic}
hold for every $n$ at no extra cost, so they are stated in that
generality. Section~\ref{P-section:n2} is new. It proves that the part
of the algebra supported on $SL(2,\cO)$ transfers exactly when
$p\nmid n$ (Proposition~\ref{P-prop:Ksub}), and computes the signs for
$n=2$ (Lemma~\ref{P-lemma:signs}). The attempt to prove the main
theorem of [Jo98] for $n=2$ led instead to a proof that it fails there
too, for constants of absolute value one
(Theorem~\ref{P-thm:unimodular}). That proof uses the new
Lemma~\ref{P-lemma:nonsplit}, that the cover does not split.
\item \emph{Extend to $n>2$ and show by explicit computations and a
counterexample which analogs fail.} Done as Part~II. The torus
commutator, the trace identity and the Iwahori-level supports are the
analogs that fail. The running example $\mathbb{Q}_7$ now serves both
parts, with $n=2$ and $n=3$.
\item \emph{First sentence of Section~8 unclear.} Rewritten as the
opening paragraph of Section~\ref{P-section:reps}.
\item \emph{Missing introductory paragraph in Section~6.3.} Added. The
same check found subsections 3.1--3.4, 6.1 and 6.2 without an
orienting paragraph. These were added too, as were the openings of the
new sections and parts.
\item \emph{Is Theorem~9.3 stated as in [Sa88]?} No. It follows the
description of [Sa88] in [Ka21], and [Sa88] itself was not consulted.
The heading now says so. The sentence on the groups of $F$-points is
our own gloss and is now Remark~\ref{P-rem:savin}.
\end{enumerate}

\section{Changes from Draft 3 to Draft 4}
\label{sec:d3d4}

\begin{enumerate}[leftmargin=1.5em]
\item The paper is divided into Part~I ($n=2$) and Part~II ($n\geq 3$).
The abstract, introduction, organization paragraph, literature review
and summary are rewritten accordingly.
\item New results: Lemma~\ref{P-lemma:nonsplit} with
Example~\ref{P-ex:nonsplit}, Section~\ref{P-section:n2}
(Proposition~\ref{P-prop:Ksub}, Lemma~\ref{P-lemma:signs},
Theorem~\ref{P-thm:unimodular}, Corollary~\ref{P-cor:unimodular} and
three examples). Remark~\ref{P-vrem:kubota} now also cites Kubota for
the splitting over $SL(2,\cO)$.
\item Theorem~\ref{P-thrm:main}, Corollary~\ref{P-corollary:main},
Example~\ref{P-ex:q7} and the remark after it moved from Section~7 to
the new Section~\ref{P-section:torus}.
\item Question~\ref{P-q:n2} rewritten in light of
Theorem~\ref{P-thm:unimodular}. The example over $\mathbb{Q}_2$ was
removed, since Part~I assumes $p$ odd.
\item The running example is $\mathbb{Q}_7$ with $n=2$ and $n=3$.
Example~\ref{P-ex:hilbert} gains the values for $n=2$.
\item Items 3--5 of the previous section. Remark~\ref{P-vrem:savin}
updated.
\item The notation table gains rows for $s$ and $w(x)$. The
acknowledgments are updated.
\end{enumerate}

\section{Changes from Draft 2 to Draft 3}
\label{sec:d2d3}

\begin{enumerate}[leftmargin=1.5em]
\item The footnote on the first page now describes the paper as a
significant revision of [Jo98], correcting, changing and expanding on
various topics, at David Joyner's request, since the title has changed.
\item The table of contents was reformatted. Subsections are indented
under their sections, the unnumbered headings inside the introduction
and the summary are omitted, and the duplicate entries for the glossary
and the acknowledgments are removed.
\item The glossary was rewritten, with an entry for Kubota's cocycle
added.
\item Question~\ref{P-q:n2} was sharpened.
\item Clarifications in the proofs of Lemma~\ref{P-lemma:cover}(b) and
Lemma~\ref{P-lemma:subgroups}(e), in the paragraph on $\exp$ and $\log$
before Definition~\ref{P-def:deep}, and in the comparison of $s_1$ with
[HM85].
\item The subtitle, acknowledgments and this log were updated for
Draft~3.
\end{enumerate}

\section{Changes from Draft 1 to Draft 2}
\label{sec:d1d2}

\begin{enumerate}[leftmargin=1.5em]
\item Title changed at David Joyner's request to ``On Hecke algebras of
the $p$-adic metaplectic cover $\overline{SL_2}$''. The subtitle and
date line follow the style guide.
\item Claude added as co-author. David Joyner's address changed to
\texttt{wdjoyner@gmail.com}. The 1998 institutional address removed.
\item The note on the original publication moved to a
\texttt{\textbackslash thanks} footnote. The subsection comparing the
paper with [Jo98] removed from the introduction. The revision log moved
from an appendix to this file.
\item New sections: literature review, notation table, summary of
results and open questions. The background section now has an example
in each subsection, using the running example $F=\mathbb{Q}_7$, $n=3$.
The glossary moved from an appendix to the end, before the
acknowledgments, and gained entries.
\item\label{it:coverb} \emph{Proof correction.} In the proof of
Lemma~\ref{P-lemma:cover}(b), Draft~1 computed $(a,a-1)_n$ as
$(a,-1)_n^{-1}(a,1-a)_n$. The correct value is $(a,-1)_n$, and
$\beta=(a,-1)_n^{-2}=1$. The statement is unchanged.
\item In the proof of Lemma~\ref{P-lemma:genuine}, the deduction that
$g_1g_2$ is admissible is now spelled out.
\item In the proof of A(i) in Proposition~\ref{P-prop:main}, the
support computation now includes the left factor $B_k$.
\item Lemma~\ref{P-lemma:length} and relation C(iv) are stated for
$\pm W_a$, which the generation argument uses.
\item Savin's theorem: the hypothesis is now $\gcd(p,n)=1$, as in the
secondary sources, and $G/Z_n$ is described as a quotient of algebraic
groups.
\item Citation keys changed to author initials plus year. Unverified
citations flagged as Remarks~\ref{P-vrem:kubota},
\ref{P-vrem:borel} and \ref{P-vrem:savin}.
\item Operators declared with \texttt{\textbackslash
DeclareMathOperator}. Equations numbered within sections. Semicolons
removed from the prose.
\item References added: [BK98], [GS12], [Ka21], [LS10], [Sa04],
[TW18], [Wa24]. Borel's page range corrected to 233--259.
\end{enumerate}

\section{Sources}
\label{sec:sources}

\subsection*{Not consulted in full text}
None of the cited works was read in full during this revision. In
particular the following statements rest on the 1998 text or on memory
(level~3): Kubota's cocycle property and local splitting
(Remark~\ref{P-vrem:kubota}), the Howe--Moy references in
Lemma~\ref{P-lemma:general} and in the remark after
Example~\ref{P-ex:t0}, and the section numbers of Borel and of
Bushnell--Henniart in Proposition~\ref{P-prop:borel}. The statement of
Savin's theorem rests on secondary sources
(Remark~\ref{P-vrem:savin}). The summaries in the literature review of
[Ge76], [RS77], [Wa80], [Fl80], [KP84], [FK86] and [Sc98] are level~3.
The summaries of [GS12], [TW18], [Ka21], [LS10] and [Wa24] rest on
their abstracts or introductions (level~2 for what is stated).

\subsection*{The summary of [Sa88] supplied by Carbon}
For Draft~4 Carbon supplied a written summary of [Sa88]. It is an AI
summary, not the source, and it was used only for orientation. Two of
its claims were checked. It agrees with [Ka21] that the target group
is $G/Z_n$. Its description of the support for $SL(2)$, the translations
$t_0^m$ with $n\mid 2m$, agrees with Proposition~\ref{P-prop:level0},
which is proved in the paper (level~1). Its other claims, including the
proof strategy through Bernstein's presentation, the Gauss sum
computation and the remarks on later refinements, were not checked and
are not used. The EuDML record of [Sa88] lists keywords that refer to
split reductive groups and to generators and relations, consistent
with the summary.

\subsection*{Bibliographic data checked online}
Volume and page data were checked for [Bo76], [Sa88], [Sa04], [Wa80],
[BK98], [GS12], [TW18], [Ka21], [LS10]. For [Bo76] the article itself
gives pages 233--259, while some indexes give 233--260. We use the
article. Whether Schultz's thesis was completed in 1998 was not checked.

\section{Style guide departures}

\begin{itemize}[leftmargin=1.5em]
\item Times is provided by \texttt{mathptmx}, because \texttt{newtx} is
not installed in the build environment. The guide allows this
fallback.
\item The title is the new one requested by David Joyner, not the
original title (guide 6.2).
\item The central result of [Jo98] is replaced, not preserved (guide
6.1), because it is false for $n\geq 3$. This was David Joyner's
decision.
\item The paper is divided into two parts, at David Joyner's request.
\item The introduction states that the main assertion of [Jo98] is
false. The style guide asks for revision history to stay in the log
(guide 6.4). This sentence is kept in the body at David Joyner's
request, because explaining that failure is the purpose of the
revision.
\item One harmless underfull vertical box remains. There are no
overfull boxes and no undefined references.
\end{itemize}

\end{document}
